% Part: first-order-logic
% Chapter: completeness
% Section: identity

\documentclass[../../../include/open-logic-section]{subfiles}

\begin{document}

\olfileid{fol}{com}{ide}
\olsection{অভিন্নতা}

\begin{explain}
আগের অংশে দেওয়া পদ-মডেলের নির্মাণটি~$\eq$-বিহীন সেট~$\Gamma$-র জন্য
প্রথম-ক্রমের যুক্তিবিদ্যার পূর্ণতা প্রতিষ্ঠা করতে যথেষ্ট। পদ-মডেলটি
প্রতিটি $!A \in \Gamma^*$-কে পরিতৃপ্ত করে, যদি বাক্যটিতে~$\eq$ না থাকে
(এবং তাই সব $!A \in \Gamma$-কেও)। কিন্তু~$\eq$ উপস্থিত থাকলে এটি কাজ
করে না। কারণ,
তখন $\Gamma^*$-তে কোনো !!a{sentence}~$\eq[t][t']$ থাকতে পারে, অথচ
পদ-মডেলে যেকোনো পদের মান সেই পদ নিজেই। তাই $t$ ও $t'$ ভিন্ন পদ হলে
পদ-মডেলে তাদের মান—যথাক্রমে $t$ ও $t'$—ভিন্ন; ফলে $\eq[t][t']$ মিথ্যা।
তবে “ভাগকরণ” নামে পরিচিত একটি নির্মাণ ব্যবহার করে এটি ঠিক করা যায়।
\end{explain}

\begin{defn}
  ধরা যাক $\Gamma^*$ হলো~$\Lang L$-এর বাক্যের একটি সঙ্গত ও !!{complete}
  সেট। $\Lang L$-এর বদ্ধ পদগুলির সেটে $\approx$ সম্পর্কটি এভাবে
  সংজ্ঞায়িত করি:
  \[
  t \approx t' \text{\quad যদি এবং কেবল যদি \quad} \eq[t][t'] \in \Gamma^*
  \]
\end{defn}

\begin{prop}
\ollabel{prop:approx-equiv}
$\approx$ সম্পর্কটির নিচের ধর্মগুলি আছে:
\begin{enumerate}
\item $\approx$ প্রতিবিম্ব।
\item $\approx$ প্রতিসম।
\item $\approx$ পরিযায়ী।
\item $t \approx t'$, $f$ একটি !!a{function}, এবং $t_1$, \dots,
  $t_{i-1}$, $t_{i+1}$, \dots, $t_n$ বদ্ধ পদ হলে
\[
\Atom{f}{t_1,\dots, t_{i-1}, t, t_{i+1}, \dots, t_n} \approx
\Atom{f}{t_1,\dots, t_{i-1}, t', t_{i+1}, \dots, t_n}.
\]
\item $t \approx t'$, $R$ একটি !!a{predicate}, এবং $t_1$, \dots,
  $t_{i-1}$, $t_{i+1}$, \dots, $t_n$ বদ্ধ পদ হলে
\begin{multline*}
\Atom{R}{t_1,\dots, t_{i-1}, t, t_{i+1}, \dots, t_n} \in \Gamma^* \text{ যদি এবং কেবল যদি } \\
\Atom{R}{t_1,\dots, t_{i-1}, t', t_{i+1}, \dots, t_n} \in \Gamma^*.
\end{multline*}
\end{enumerate}
\end{prop}

\begin{proof}
$\Gamma^*$ সঙ্গত ও !!{complete} বলে $\eq[t][t'] \in \Gamma^*$ যদি এবং
কেবল যদি $\Gamma^* \Proves \eq[t][t']$। তাই নিচের কথাগুলি দেখালেই যথেষ্ট:
\begin{enumerate}
\item প্রতিটি বদ্ধ পদ~$t$-এর জন্য $\Gamma^* \Proves \eq[t][t]$।
\item $\Gamma^* \Proves \eq[t][t']$ হলে $\Gamma^* \Proves \eq[t'][t]$।
\item $\Gamma^* \Proves \eq[t][t']$ এবং $\Gamma^* \Proves
  \eq[t'][t'']$ হলে $\Gamma^* \Proves \eq[t][t'']$।
\item $\Gamma^* \Proves \eq[t][t']$ হলে
\[
\Gamma^* \Proves
\eq[\Atom{f}{t_1,\dots,t_{i-1},t,t_{i+1},\dots,t_n}][\Atom{f}{t_1,\dots,t_{i-1},t',t_{i+1},\dots,t_n}]
\]
প্রতিটি $n$-স্থানীয় !!{function}~$f$ এবং বদ্ধ পদ $t_1$, \dots,
$t_{i-1}$, $t_{i+1}$, \dots,~$t_n$-এর জন্য।% BN-SRC-079 TARGET-CORRECTION-BEGIN
\par\smallskip\noindent\textnormal{\small\textbf{উৎস-সংশোধন (BN-SRC-079):} হিমায়িত ইংরেজি প্রদর্শিত অপেক্ষক-পদটির প্রথম পাশে $t_{i+1}$-এর পরে দুটি কমা আছে। বাংলা লক্ষ্যে অনিচ্ছাকৃত অতিরিক্ত কমাটি বাদ দিয়ে দুই পাশের যুক্তি-অনুক্রম সমান্তরাল রাখা হয়েছে।} % BN-SRC-079 TARGET-CORRECTION-END
\item $\Gamma^* \Proves \eq[t][t']$ এবং
$\Gamma^* \Proves
\Atom{R}{t_1,\dots,t_{i-1},t,t_{i+1},\dots,t_n}$ হলে
$\Gamma^* \Proves \Atom{R}{t_1,\dots,t_{i-1},t',t_{i+1},\dots,t_n}$,
প্রতিটি $n$-স্থানীয় !!{predicate}~$R$ এবং বদ্ধ পদ $t_1$, \dots,
$t_{i-1}$, $t_{i+1}$, \dots,~$t_n$-এর জন্য।
\end{enumerate}
\end{proof}

\begin{prob}
\olref[fol][com][ide]{prop:approx-equiv}-এর প্রমাণ সম্পূর্ণ করো।
\end{prob}

\begin{defn}
ধরা যাক $\Gamma^*$ কোনো ভাষা~$\Lang L$-এর একটি সঙ্গত ও !!{complete}
সেট, $t$ একটি বদ্ধ পদ, এবং $\approx$ আগের সংজ্ঞার মতো। তাহলে:
\[
\equivrep{t}{\approx} = \Setabs{t'}{t'\in \Trm[L], t \approx t'}
\]
এবং $\equivclass{\Trm[L]}{\approx} = \Setabs{\equivrep{t}{\approx}}{t \in \Trm[L]}$।
\end{defn}

\begin{defn}
\ollabel{defn:term-model-factor}
\olref[mod]{defn:termmodel} থেকে~$\Gamma^*$-র পদ-মডেলকে $\Struct M =
\Struct M(\Gamma^*)$ ধরি। তাহলে $\Struct{\equivclass{M}{\approx}}$ হলো
নিচের !!{structure}:
\begin{enumerate}
\item $\Domain{\equivclass{M}{\approx}} = \equivclass{\Trm[L]}{\approx}$।
\item $\Assign{c}{\equivclass{M}{\approx}} = \equivrep{c}{\approx}$
\item $\Assign{f}{\equivclass{M}{\approx}}(\equivrep{t_1}{\approx}, \dots,
  \equivrep{t_n}{\approx}) = \equivrep{\Atom{f}{t_1,\dots, t_n}}{\approx}$
\item $\tuple{\equivrep{t_1}{\approx}, \dots, \equivrep{t_n}{\approx}} \in
  \Assign{R}{\equivclass{M}{\approx}}$ যদি এবং কেবল যদি
  $\Sat{M}{\Atom{R}{t_1,\dots, t_n}}$; অর্থাৎ, যদি এবং কেবল যদি
  $\Atom{R}{t_1,\dots, t_n} \in \Gamma^*$।
\end{enumerate}
\end{defn}

\begin{explain}
লক্ষ করো, $\equivclass{\Trm[L]}{\approx}$-এর উপাদানগুলিকে
$\equivrep{t}{\approx}$ হিসেবে উল্লেখ করে—অর্থাৎ
\emph{প্রতিনিধি}~$t \in \equivrep{t}{\approx}$-এর মাধ্যমে—আমরা
$\Assign{f}{\equivclass{M}{\approx}}$ ও
$\Assign{R}{\equivclass{M}{\approx}}$ সংজ্ঞায়িত করেছি। নিশ্চিত করতে হবে
যে সংজ্ঞাগুলি এই প্রতিনিধিদের নির্বাচনের ওপর নির্ভর করে না; অর্থাৎ একই
তুল্যতা শ্রেণী নির্ধারণকারী অন্য কোনো নির্বাচন~$t'$-এর ক্ষেত্রে
($\equivrep{t}{\approx} = \equivrep{t'}{\approx}$) সংজ্ঞাগুলি একই ফল
দেয়। যেমন, $R$ একটি এক-স্থানীয় !!{predicate} হলে সংজ্ঞার শেষ দফা বলে
$\equivrep{t}{\approx} \in \Assign{R}{\equivclass{M}{\approx}}$ যদি এবং
কেবল যদি $\Sat{M}{\Atom{R}{t}}$। $t \approx t'$ এমন অন্য কোনো পদ~$t'$-এর
জন্য যদি $\Sat/{M}{\Atom{R}{t'}}$ হয়, তবে সংজ্ঞাটি দাবি করত
$\equivrep{t'}{\approx} \notin \Assign{R}{\equivclass{M}{\approx}}$।
$t \approx t'$ হলে $\equivrep{t}{\approx} =
\equivrep{t'}{\approx}$; কিন্তু একই সঙ্গে $\equivrep{t}{\approx} \in
\Assign{R}{\equivclass{M}{\approx}}$ ও $\equivrep{t}{\approx} \notin
\Assign{R}{\equivclass{M}{\approx}}$ হতে পারে না। তবে
\olref{prop:approx-equiv} নিশ্চিত করে যে এমনটি ঘটতে পারে না।% BN-SRC-080 TARGET-CORRECTION-BEGIN
\par\smallskip\noindent\textnormal{\small\textbf{উৎস-সংশোধন (BN-SRC-080):} হিমায়িত ইংরেজি প্রতিনিধি-স্বাধীনতার উদাহরণে “অন্য পদ~$t'$” প্রবর্তন করেও মিথ্যা হওয়া পরমাণু বাক্যে আবার $R(t)$ লেখে। নির্বাচনের ওপর নির্ভরতার সম্ভাব্য বিরোধটি তুলনা করতে সেখানে $R(t')$ আবশ্যক; বাংলা লক্ষ্যে সেই প্রতিনিধিটি পুনঃস্থাপন করা হয়েছে।} % BN-SRC-080 TARGET-CORRECTION-END
\end{explain}

\begin{prop}
$\Struct{\equivclass{M}{\approx}}$ সুসংজ্ঞায়িত। অর্থাৎ $t_1$,
  \dots, $t_n$, $t_1'$, \dots, $t_n'$ বদ্ধ পদ এবং প্রতিটি সূচকের জন্য
  $t_i \approx t_i'$ হলে
\begin{enumerate}
\item $\equivrep{\Atom{f}{t_1,\dots, t_n}}{\approx} =
    \equivrep{\Atom{f}{t_1',\dots, t_n'}}{\approx}$; অর্থাৎ,
  \[
  \Atom{f}{t_1,\dots, t_n} \approx \Atom{f}{t_1',\dots, t_n'}
  \]
  এবং
\item $\Sat{M}{\Atom{R}{t_1,\dots, t_n}}$ যদি এবং কেবল যদি
  $\Sat{M}{\Atom{R}{t_1',\dots, t_n'}}$; অর্থাৎ,
  \[
    \Atom{R}{t_1,\dots, t_n} \in \Gamma^* \text{ যদি এবং কেবল যদি }
    \Atom{R}{t_1',\dots, t_n'} \in \Gamma^*.
  \]
\end{enumerate}
\end{prop}

\begin{proof}
$n$-এর ওপর আরোহের সাহায্যে \olref{prop:approx-equiv} থেকে অনুসৃত হয়।
\end{proof}

পদ-মডেলের ক্ষেত্রের মতোই, সত্যতা-সহায়ক উপপাদ্য প্রমাণের আগে নিচের
সহায়ক উপপাদ্যটি দরকার।

\begin{lem}
\ollabel{lem:val-in-termmodel-factored} $\Struct M = \Struct M
 (\Gamma^*)$ হলে প্রতিটি বদ্ধ পদের জন্য
$\Value{t}{\equivclass{M}{\approx}} = \equivrep{t}{\approx}$।% BN-SRC-081 TARGET-CORRECTION-BEGIN
\par\smallskip\noindent\textnormal{\small\textbf{উৎস-স্পষ্টীকরণ (BN-SRC-081):} হিমায়িত ইংরেজি ভাগফল-পদ-মডেলের মান-সহায়ক উপপাদ্যে $t$-কে অবাধ রাখা হয়েছে, অথচ চল-মূল্যায়ন ছাড়া প্রদত্ত মানায়ন এবং সংশ্লিষ্ট \olref[mod]{lem:val-in-termmodel} কেবল বদ্ধ পদের জন্য। বাংলা লক্ষ্যে সিদ্ধান্তে সেই পরিসর স্পষ্ট করা হয়েছে।} % BN-SRC-081 TARGET-CORRECTION-END
\end{lem}
\begin{proof}
প্রমাণটি \olref[mod]{lem:val-in-termmodel}-এর প্রমাণের অনুরূপ।
\end{proof}

\begin{prob}
\olref[fol][com][ide]{lem:val-in-termmodel-factored}-এর প্রমাণ সম্পূর্ণ করো।
\end{prob}

\begin{lem}
\ollabel{lem:truth}
সব বাক্য~$!A$-এর জন্য $\Sat{\equivclass{M}{\approx}}{!A}$ যদি এবং কেবল
যদি $!A \in \Gamma^*$।
\end{lem}

\begin{proof}
\olref[mod]{lem:truth}-এর প্রমাণের মতোই, $!A$-এর ওপর আরোহের সাহায্যে।
শুধু $!A \ident \eq[t][t']$ ক্ষেত্রটির জন্য অতিরিক্ত মনোযোগ দরকার।
\begin{align*}
\Sat{\equivclass{M}{\approx}}{\eq[t][t']} & \text{ যদি এবং কেবল যদি } \equivrep{t}{\approx} = \equivrep{t'}{\approx}
\text{ ($\Struct{\equivclass{M}{\approx}}$-এর সংজ্ঞা অনুসারে)}\\
& \text{ যদি এবং কেবল যদি } t \approx t' \text{ ($\equivrep{t}{\approx}$-এর সংজ্ঞা অনুসারে)}\\
& \text{ যদি এবং কেবল যদি } \eq[t][t'] \in \Gamma^* \text{ ($\approx$-এর সংজ্ঞা অনুসারে)।}
\end{align*}
\end{proof}

\begin{digress}
লক্ষ করো, $\Struct{M(\Gamma^*)}$ সব সময় !!{enumerable} ও অসীম হলেও
$\Struct{\equivclass{M}{\approx}}$ সসীম হতে পারে; কারণ
$\equivrep{t}{\approx}$ শ্রেণী কেবল সসীমসংখ্যকও হতে পারে। এমনটিই
প্রত্যাশিত, কারণ~$\Gamma$-তে এমন !!{sentence}s থাকতে পারে যেগুলি দাবি
করে যে তাদের সত্যকারী যেকোনো !!{structure} সসীম। যেমন,
$\lforall[x][\lforall[y][x = y]]$ একটি সঙ্গত !!{sentence}, কিন্তু সেটি
শুধু ঠিক একটি !!{element}-বিশিষ্ট !!a{domain}-এর !!{structure}s-তেই
পরিতৃপ্ত হয়।
\end{digress}

\end{document}
